\section{Certified commutation of closed register fragments}

This note supplies the ordinary semantic argument and an independent
certificate checker for a grammar search modulo commuting register
fragments. The underlying idea of disjoint-register trace grammars already
appears in ProveIt intake163/source66 and source67. The refinement examined
here permits shared additive outputs while forbidding interference with
tested registers. It is a sufficient condition, not a complete procedure
for deciding whether two programs commute. No historical priority claim
or Turing-machine state saving is attached to the condition itself.

\subsection{Effect condition and partial-map proof}

A fragment is a fixed structured register program over nine natural-number
registers, using increments, decrement/zero tests, structured breaks and
continues. The frozen source DAG's structural analysis verifies that a
fragment is closed: no break or continue escapes to an enclosing source
context. Exclude every fragment containing the primitive \texttt{halt}.
A remaining fragment either returns normally with a register file or runs
forever, and consequently denotes a partial function on $\mathbb N^9$.

For a fragment $P$, let $W(P)$ be the set of registers appearing in an
increment or decrement instruction, and let $D(P)$ be those appearing in
a decrement/zero-test instruction. The conservative independence condition
is
\begin{equation}\label{trace-additive-condition}
D(P)\cap W(Q)=\varnothing,
\qquad
D(Q)\cap W(P)=\varnothing.
\end{equation}
The strict control instead requires $W(P)\cap W(Q)=\varnothing$.

\paragraph{Commutation theorem.}
If two closed, halt-free fragments satisfy
\eqref{trace-additive-condition}, then their partial normal-return
functions commute: $P;Q$ and $Q;P$ have the same domain and the same
output register file on that domain.

\paragraph{Proof.}
Every control decision in $P$ depends only on a register in $D(P)$.
Initial values of all other registers cannot change its instruction path
or termination. More precisely, keeping the initial $D(P)$ coordinates
fixed keeps the entire control trace fixed. On a terminating trace, each
coordinate outside $D(P)$ is unchanged except for a fixed finite number
of increments, whose count depends only on those initial $D(P)$ values.
The analogous statements hold for $Q$.

By the first condition, $Q$ never modifies an input determining $P$'s
trace; by the second, $P$ never modifies one determining $Q$'s trace.
Thus both fragments terminate in either order iff each terminates on the
initial register file. On those initial files, a register in $D(P)$ is
untouched by $Q$, and a register in $D(Q)$ is untouched by $P$. On every
remaining register, each fragment contributes only its fixed nonnegative
increment count, and the two additions commute. This proves equality of
the outputs. If either fragment diverges, its trace is unaffected by the
other and neither composition returns. No totality assumption is needed.

\paragraph{Why the boundaries matter.}
A primitive halt cannot be included merely because its register effect
is empty: a halting fragment followed by a divergent fragment would halt
in one order and diverge in the other. Likewise a break that escapes the
fragment is not a partial function from register files to normal returns;
its effect depends on an enclosing continuation. The checker excludes
halts and verifies closure before applying any commutation rule. It
permits rearrangement of complete source atoms, not arbitrary pieces of
the tape decoder or return implementation.

The structural closure analysis tracks a pair of outward break and
continue bounds. Its leaves assign bounds $(1,0)$ to a decrement,
$(n,0)$ to \texttt{brk} $n$, $(0,n)$ to \texttt{cont} $n$, and $(0,0)$
to increment or halt. For a source composition node with child bounds
$a,b$, the pair is
\[
\bigl(\max(a_0-1,b_0,0),\;\max(a_1,b_1-1,0)\bigr).
\]
All forty-four current semantic atoms have bound $(0,0)$. Atoms
\texttt{J} and \texttt{n} are excluded because they contain halt.
The proof uses the existing semantics of this structural source analysis;
the emitted effect data is not a replacement for the source compiler's
correctness argument.

\subsection{Canonical words and independently checkable certificates}

Use a symmetric, irreflexive independence relation on the semantic atom
alphabet. Two words are trace-equivalent when adjacent independent
letters can be swapped. The partial-map theorem makes every such swap
sound in a sequence of complete closed fragments, including a sequence
whose later continuation eventually halts.

The normalizer uses the lexicographically least topological order of the
word's dependence DAG. Its vertices are occurrences. For every earlier
occurrence of a letter dependent on a later letter there is a precedence
constraint. It suffices to store an edge from the last earlier occurrence
of each dependent letter: repeated occurrences of one letter are already
chained, so transitive closure implies all earlier constraints. A priority
queue of currently minimal occurrences produces a unique canonical word.
Ordinary adjacent-inversion sorting is not a valid substitute: a partial
commutation relation need not make such a rewriting strategy confluent.

The normal form $\operatorname{NF}$ satisfies
\[
\operatorname{NF}(\operatorname{NF}(u)\operatorname{NF}(v))
 =\operatorname{NF}(uv),
\qquad
\operatorname{NF}(u)=\operatorname{NF}(v)
\ \Longleftrightarrow\ u\text{ and }v\text{ are trace-equivalent}.
\]
Thus an intern table in a grammar search may use
$\operatorname{NF}(uv)$ as the key of a composition, provided every rule
continues to denote the concatenation of the two child functions.
Rotations, alternative splits and extraction of repeated subwords then
operate in this quotient monoid.

The companion \texttt{trace\_commutation.py} independently derives every
effect from the frozen AST. For each root word it outputs an explicit
list of adjacent-swap positions, together with the original and final
word. A small checker replays that list, checks independence at every
step against the derived effect relation, and verifies the final word.
This avoids trusting the searcher's normalizer merely because it reports
the same key. Further grammar-to-source and binary compiler verification
is still required for a newly emitted machine.

\subsection{Finite controls on the 245-state input grammar}

The checked input has 128 grammar rules, fourteen roots and forty-four
semantic atoms. The exact observations are:
\begin{center}
\begin{tabular}{lrr}
\toprule
Observation & Strict disjoint & Additive outputs\\
\midrule
Unordered independent atom pairs &162&211\\
Eligible adjacent pairs across roots &70&91\\
Roots whose canonical spelling changes &9&9\\
Adjacent swaps in all root certificates &154&154\\
Already duplicate grammar keys after normalization &0&0\\
\bottomrule
\end{tabular}
\end{center}
The larger independence relation gives additional legal search moves but
does not by itself collapse any of the current 128 keys. The script also
checks normal-form idempotence, the concatenation identity at every split
of every current grammar word, and the swap identity for every independent
pair. The 211 additive-output pairs agree exactly with the independently
generated relation used by the C++ search lane.

Pinned inputs, with SHA256 values, are:
\begin{itemize}
\item \texttt{closed-sequence-grammar-next22.json}:\\
\texttt{3d7798e924a9e5d56594d29e8a5c46bf1e9affc3f23fa4e6a1695499112fac19}.
\item \texttt{semantic-symbols.json}:\\
\texttt{4e77c352494a3e18d143776a1b5ffca1a2e7d234aba0c2560a41982339395fdd}.
\item \texttt{state\_reduction/route-grammar.txt}:\\
\texttt{6e52c8830dfa46b3e38a3a6d081e91316521393c7029a0b115a2f4d2e3d3398c}.
\end{itemize}
Reproduce with the source-proof input directory and routed grammar:
\begin{verbatim}
python trace_commutation.py --inputs INPUTS \
  --grammar INPUTS/state_reduction/route-grammar.txt \
  --output trace_commutation_results.json
\end{verbatim}
The results include all atom footprints and root swap certificates.
The source input and symbols belong to the earlier
\texttt{zfc245-proof-20261018} package; this note changes neither input.
